\chapter{史料与参考路径}
\section{先读哪几类史书}
\term{二十四史}提供历代王朝的官修叙事主线：从\sourcebook{史记}、\sourcebook{汉书}、\sourcebook{后汉书}、\sourcebook{三国志}，到\sourcebook{旧唐书}、\sourcebook{新唐书}、\sourcebook{宋史}、\sourcebook{元史}、\sourcebook{明史}。它们适合查人物、制度、战争和一朝的自我叙述，但也各自带着成书时代的立场。

\term{\sourcebook{资治通鉴}}把战国至五代的材料重新按年组织，是本书编年写法最重要的参照：读者能看到一件事之前发生了什么、之后又引出什么。但《通鉴》并不是不带判断的事件清单，它的取舍与司马光的政治关切同样需要被看见。

\term{补充材料}包括\sourcebook{续资治通鉴长编}、\sourcebook{建炎以来系年要录}、\sourcebook{明实录}、\sourcebook{清实录}，以及金文、简帛、地方文书和现代专题研究。越早的时代越应重视出土材料；越近的时代越要注意官方档案与私人记述之间的差别。

\section{引用原则}
书稿正式扩写时，每一条深描事件至少标明一项一手或近一手材料与一项现代研究；纪年有分歧时，优先说明分歧而非伪造精确性。正文会尽量把史料融进讲述里，而不是把读者抛进密密麻麻的引文墙。
